\begin{background}
Network intrusion detection systems (NIDS) have become an essential component of modern cybersecurity infrastructure, safeguarding computer networks against a continuously expanding range of malicious activities. As network traffic volumes grow and attack patterns become more sophisticated, traditional signature-based detection methods struggle to keep pace. Machine learning and deep learning techniques offer a viable alternative, capable of learning intricate patterns from historical network data and generalizing to identify previously unseen threats. The NSL-KDD dataset, a refined successor to the original KDD Cup 1999 benchmark, remains one of the most widely adopted resources for evaluating intrusion detection algorithms because of its standardized format, balanced class distribution, and explicit separation between training and testing partitions.

Internationally, a substantial body of research has applied diverse machine learning and deep learning models to the NSL-KDD benchmark, including decision trees, random forests, support vector machines, multi-layer perceptrons, convolutional neural networks, and long short-term memory networks. Despite the encouraging results reported in many published studies, several persistent challenges continue to limit practical deployment. Models trained on historical attack signatures frequently fail to generalize to novel attack categories absent from the training data, a difficulty commonly referred to as the zero-day or out-of-distribution detection problem. The pronounced class imbalance between normal traffic and rare attack types also produces biased classifiers and inflated accuracy metrics. Moreover, inconsistent evaluation protocols across published works make rigorous comparison of competing methods difficult.

This paper reports a systematic comparison of ten representative models on the NSL-KDD dataset, covering both traditional machine learning algorithms (decision tree, random forest, $K$-nearest neighbors, naive Bayes, logistic regression, support vector machine, and multi-layer perceptron) and deep learning architectures (one-dimensional convolutional neural network and long short-term memory network). Experiments use the standard 125{,}973 training samples and 22{,}544 test samples with 41 features and 23 attack types, evaluating both binary classification (normal versus attack) and five-class multi-class classification. All models share a unified preprocessing pipeline and an identical test partition, ensuring fair and reproducible comparison. Performance is documented through 15 evaluation figures that include accuracy, precision, recall, $F_1$-score, confusion matrices, and per-class breakdowns.

The work originates from a graduate-level network security course project undertaken by a three-member team. The team previously built the complete machine learning pipeline from data exploration and preprocessing through model training, hyperparameter tuning, and comprehensive evaluation, with all code implemented in Python using scikit-learn and PyTorch. This paper extends that course project by formalizing the experimental design, broadening model coverage, and providing a critical analysis of where deep learning genuinely helps and where simpler methods remain competitive. Rather than reporting only favorable outcomes, the paper offers a candid assessment of deep learning limitations on out-of-distribution test data and concludes with practical deployment recommendations for resource-constrained network environments.
\end{background}
